% !PRINTED 120
\weberpage{120}
Wir zerlegen $P$ wie in (4) in die $j$ Nebengruppen $B$, bezeichnen
mit $e$ alle Elemente von $Q$, und untersuchen nun die Permutation
\[
\pi_e = (B,\ Be),
\]
die durch die Elemente $e$ unter den $j$ Elementen (4) hervorgerufen
wird. Diese Permutationen bilden nach (5) und (6) gewiss eine
mit $Q$ ein- oder mehrstufig isomorphe Gruppe $\Pi$.

Ist $a_1$ irgend ein Element von $P$, so werden in dem nach
der Composition der Theile gebildeten Complex
\[
P_1 = Ra_1Q
\]
eine gewisse Anzahl der Nebengruppen $B$ enthalten sein. Diese
Anzahl finden wir, wenn wir die Bedingung aufsuchen, dass
$Ra_1 = Ra_1e$ sei, oder dass $a_1e$ in $Ra_1$ oder $e$ in $a_1^{-1}Ra_1$ enthalten
sei. Ist also
\[
Q_1 \text{ der Durchschnitt von } Q \text{ und } a_1^{-1}Ra_1
\]
vom Index $h_1$ in Bezug auf $Q$, so enthält $P_1$ genau $h_1$ und nicht
mehr von den Complexen $B$.

Ist $a_2$ ein zweites Element von $P$, so ist der Complex
\[
P_2 = Ra_2Q
\]
entweder völlig mit $P_1$ übereinstimmend, oder beide enthalten
gar kein gemeinsames Element. Denn wenn $P_1$ und $P_2$ ein
gemeinsames Element enthalten, so haben sie einen der Complexe $B$ gemein. Das ist aber nur möglich, wenn $a_2e$ in $Ra_1e'$,
also $a_2$ in $P_1$ enthalten ist, und dann ist $P_1$ und $P_2$ identisch.
Wir können also mit der Bildung der Complexe $P_1, P_2 \ldots$ fortfahren, bis die Gruppe $P$ erschöpft ist, und erhalten die Zerlegung der Gruppe $P$ nach den zwei Gruppen $Q$ und $R$:
\weq{9}{\begin{aligned}
P &= Ra_1Q + Ra_2Q + \cdots\\
&= P_1 \phantom{{}={}} + P_2 \phantom{{}={}} + \cdots
\end{aligned}}

Aus der Bemerkung, dass, wenn $a_2$ nicht in $Ra_1Q$ vorkommt,
$a_2^{-1}$ nicht in $Qa_1^{-1}R$ enthalten ist, ergiebt sich die zweite
Zerlegung:
\weq{10}{P = Qa_1^{-1}R + Qa_2^{-1}R + \cdots\footnotemark.}

Die Bestandtheile $P_1, P_2, \ldots$ von (9) enthalten je $h_1, h_2, \ldots$
Elemente der Reihe $B$, wenn $h_1, h_2, \ldots$ die Indices der Durchschnitte von $Q$ mit $a_1^{-1}Ra_1, a_2^{-1}Ra_1, \ldots$ in Bezug auf $Q$, also % [sic: Weber prints a_2^{-1}Ra_1; the conjugate pattern requires a_2^{-1}Ra_2.]
Theiler des Grades $q$ von $Q$ bedeuten.
\footnotetext{Diese Zerlegungen verdanke ich einer Mittheilung von \gesperrt{Dedekind}.}
